\section{Chapitre~\ref{sec:nora}. \nt{NORA}}

L'organisation du système \nt{NORA}, ses deux modes, le lien avec la pupille (qui ne passe pas seulement par \nt{NORA}), l'augmentation du rapport du signal au fond et le U inversé du comportement sont mesurés directement; le gain multiplicatif $g$ est un modèle; que le gain change le régime du choix et agisse comme l'inverse d'une température sont des conclusions du chapitre; le bénéfice de l'ajustement du gain est un calcul par le modèle du livre; \og \nt{NORA}, bouton de l'exploration\fg{} et \og \nt{NORA}, interruption\fg{} sont des hypothèses.

{\footnotesize
\setlength{\tabcolsep}{3pt}\renewcommand{\arraystretch}{1.15}
\begin{longtable}{>{\raggedright}p{0.5cm}>{\raggedright}p{4.0cm}>{\raggedright}p{5.6cm}>{\raggedright}p{2.3cm}>{\raggedright\arraybackslash}p{1.7cm}}
\toprule
N\textsuperscript{o} & Affirmation & Expérience et ce qui est mesuré & Source & Statut \\
\midrule\endfirsthead
\toprule
N\textsuperscript{o} & Affirmation & Expérience et ce qui est mesuré & Source & Statut \\
\midrule\endhead
1 & L'homme possède quelques dizaines de milliers de neurones \nt{NORA}, presque tous en un seul groupe & Coupes sériées du tronc cérébral humain, marquage de la tyrosine hydroxylase: $53\,900$ cellules dans le locus coeruleus et encore $6\,260$ dans un sous-noyau voisin. & \cite{Baker1989LC} & mesuré \\
2 & Les sorties se répartissent dans presque tout le cerveau, mais les parties du groupe desservent en partie des régions différentes & Marquage viro-génétique chez la souris: les neurones \nt{NORA} du locus coeruleus reçoivent des entrées de vastes régions du cerveau et envoient des sorties vers de vastes régions du cerveau et de la moelle épinière. Chez le rat: le groupe qui projette vers l'amygdale renforce l'apprentissage de la peur, et un groupe distinct projetant vers le cortex préfrontal renforce son extinction. & \cite{Schwarz2015LC, Uematsu2017LC, Poe2020} & mesuré \\
3 & Mode tonique: impulsions régulières, d'une fraction d'impulsion à quelques impulsions par seconde, fréquence qui varie avec l'éveil & Rats, enregistrement de neurones du locus coeruleus pendant le cycle veille-sommeil: la fréquence est maximale à l'éveil, plus basse en sommeil lent, presque nulle en sommeil paradoxal; les variations de fréquence précèdent le changement d'état. & \cite{AstonJonesBloom1981} & mesuré \\
4 & Mode phasique: courte bouffée sur un événement important pour la tâche, à peu près au moment de la décision & Singes, tâche de détection d'une cible rare: tous les neurones du locus coeruleus répondent par une bouffée à la cible seulement, $\sim90\ms$ après elle et $\sim200\ms$ avant la réponse manuelle; quand la performance est mauvaise, les bouffées sont plus faibles. Dans une tâche de choix entre deux options, la bouffée est plus étroitement liée à la réponse qu'au stimulus, et elle est présente pour les réponses justes comme pour les erreurs. & \cite{AstonJones1994vigilance, Clayton2004LC} & mesuré \\
5 & Le diamètre de la pupille est un point de mesure imprécis de \nt{NORA}: il suit \nt{NORA}, mais aussi \nt{ACET} et d'autres régions & Singes: les impulsions du locus coeruleus, jusqu'aux impulsions isolées d'une seule cellule, prédisent la dilatation de la pupille; un lien semblable, mais plus tardif, existe dans les colliculus et le cortex cingulaire. Souris: les fluctuations de la pupille suivent l'activité des sorties noradrénergiques comme cholinergiques dans le cortex. & \cite{Joshi2016, Reimer2016} & mesuré \\
6 & La noradrénaline supprime l'activité de fond plus que l'activité évoquée; le gain multiplicatif~\eqref{eq:gain} est un modèle qui rend compte de cet effet & Singes éveillés, cortex auditif, iontophorèse de noradrénaline: l'activité de fond des neurones est plus supprimée que la réponse aux vocalisations des congénères; le rapport signal sur bruit augmente. La sigmoïde multiplicative~\eqref{eq:gain} est reprise d'un modèle de réseau; la multiplication littérale de la pente n'est pas mesurée. & \cite{Foote1975NE, ServanSchreiber1990} & mesuré; $g$: modèle \\
7 & Le gain n'augmente $d'$ que contre le bruit postérieur à l'amplificateur~\eqref{eq:gainsnr} (proposition~\ref{prop:gainnoise}) & Démonstration dans le chapitre; dans un modèle de réseau, le gain améliore la détection du signal. Aucune expérience ne sépare le bruit avant et après l'amplificateur pour tester cette distinction. & démonstration dans le chapitre; \cite{ServanSchreiber1990} & théorie \\
8 & La frontière $g\beta=1$ fait passer le réseau de choix de \og réfléchit\fg{} à \og a décidé\fg{}~\eqref{eq:wtagain} (proposition~\ref{prop:gainwta}, fig.~\ref{fig:pitchfork}) & Analyse linéaire de deux groupes qui s'inhibent mutuellement; démonstration dans le chapitre. Ce seuil n'a pas été mesuré directement dans le cerveau. & démonstration dans le chapitre & théorie \\
9 & Le gain est l'inverse de la température du choix~\eqref{eq:softmax} & Avec un bruit logistique après l'amplificateur, la probabilité de choix est un softmax avec $T=\sigma/g$; démonstration dans le chapitre. & démonstration dans le chapitre & théorie \\
10 & \nt{NORA} fixe combien chercher: forte activité tonique, petit $g$ effectif (exploration); bouffée phasique au moment de la décision, grand $g$ (décision) & Humains: avant un choix au hasard, la pupille de base est plus large qu'avant le choix de la meilleure option connue; les personnes à pupille plus large explorent davantage. Rats: le passage au mode de choix aléatoire est commandé par l'entrée \nt{NORA} dans le cortex cingulaire antérieur. Que la bouffée augmente le $g$ effectif n'est pas mesuré directement. & \cite{JepmaNieuwenhuis2011, Tervo2014stochastic, AstonJones2005} & hypothèse \\
11 & La récompense dépend de $g$ comme un U inversé; le meilleur $g$ est plus petit dans un monde qui change vite (proposition~\ref{prop:explore}, fig.~\ref{fig:invertedU}) & \texttt{models/nora\_explore.py}, $60$ simulations de $4000$ essais. Changement tous les $1000$ essais: meilleur $g=16$, part $0.976$; tous les $20$: $g=8$, part $0.886$; pour $g=0.5$, $0.77$. Le recalcul concorde avec le chapitre. & \texttt{models/\allowbreak nora\_\allowbreak explore.py} & modèle \\
12 & U inversé dans le comportement: meilleur résultat avec une activité tonique modérée et des bouffées nettes & Singes, même tâche qu'à la ligne~4: pendant les périodes de bonne performance, la fréquence tonique est modérée et les bouffées sur la cible sont nettes; avec une fréquence tonique élevée, il n'y a pas de bouffées et les fausses réponses sont plus nombreuses. Les variations de l'activité tonique et évoquée suivent les fluctuations de la qualité de la performance. & \cite{Usher1999LC, AstonJones2005} & mesuré \\
13 & \nt{NORA} signale l'incertitude inattendue, \nt{ACET} l'incertitude attendue & Théorie de l'inférence bayésienne à deux sortes d'incertitude; les travaux cités ne contiennent pas d'expérience directe séparant ces signaux selon le neurotransmetteur. & \cite{YuDayan2005} & hypothèse \\
14 & Après un changement brusque, les humains apprennent plus vite, et la pupille se dilate & Humains, tâche de prédiction avec changements soudains: les variations brèves de la pupille suivent l'estimation de la probabilité d'un changement et, avec la pupille de base, prédisent à quel point les nouvelles données modifient l'estimation (vitesse d'apprentissage); une variation de la pupille sans lien avec la lumière ni avec la tâche modifie cette vitesse. Que la pupille indique ici précisément \nt{NORA} est une conclusion tirée des données indirectes de la ligne~5. & \cite{Nassar2012} & mesuré \\
15 & La bouffée phasique fonctionne comme une interruption: le réseau remet son état à zéro & Aucune expérience directe ne montre une bouffée \nt{NORA} remettant à zéro l'état du réseau; seules les bouffées sur les événements importants sont mesurées (ligne~4). & \cite{BouretSara2005} & hypothèse \\
16 & L'ajustement de $g$ ou de la vitesse d'apprentissage par la surprise ne fait guère mieux que les meilleurs réglages constants & \texttt{models/nora\_explore.py}, monde à époques de $1000$ essais (changement tous les $1000$ et tous les $20$ essais). Meilleur $g$ constant, $g=8$: $0.925$; ajustement de $g$: jusqu'à $0.927$; ajustement de la vitesse d'apprentissage: jusqu'à $0.926$; vitesse constante $0.4$: $0.928$. Le recalcul concorde avec le chapitre. & \texttt{models/\allowbreak nora\_\allowbreak explore.py} & modèle \\
\bottomrule
\end{longtable}}
